\section{Reconstruction intégrale du registre et isométrie d’incidence}
\label{rev4:ki:section}

La reconstruction de la section~\ref{sec:memory} permet de conserver dans
une même composition la géométrie positive, l’incidence exceptionnelle
et la direction de déformation. L’information initiale de cette
composition est le registre produit par l’état enrichi
APP--TRIT--TPK : les deux feuillets arithmétiques conservent leurs opérations ;
l’orientation tritique détermine le régime ; le transport, la
mise à jour et le prolongement nonadique conservent les différences
entre positions et la composante uniforme. Les réalisations linéaires
qui suivent agissent sur ce registre déjà déterminé. Leur fonction est
de transporter l’information et d’expliciter ses invariants.

La distinction fondamentale concerne le domaine conservé. Le
registre \(K\) contient douze coordonnées ordonnées ; la tétrade
\(B_K\) conserve quatre positions sélectionnées par l’incidence ;
une direction géométrique conserve une droite dans un espace de
représentation. Chaque opération a une image et une fibre propres.
La construction suivante conserve d’abord les douze coordonnées et
dérive ensuite les réductions qui interviennent dans la géométrie. Ainsi,
une reconstruction exacte est disponible avant d’effectuer une
extraction de support ou une réduction dimensionnelle.

\subsection{Forme de Gram du plan dodécaphasique}
\label{rev4:ki:gram-subsection}

Nous adoptons la numérotation de un à douze et notons
\(\Omega=\{1,\ldots,12\}\) l’ensemble des positions. Le
codage Paley--Witt de l’état enrichi détermine l’ensemble
\(\mathcal H\) de ses \(132\) hexades. Une hexade est un support
de six positions d’un vecteur de poids six du code ternaire
étendu. La construction explicite du code et la démonstration du fait
que chaque quintuplet appartient à une unique hexade sont développées dans
le théorème~\ref{thm:lf-code} et la
proposition~\ref{prop:lf-witt}. Nous utilisons ici précisément cette
incidence, avec les positions transportées dans le même ordre que
les coordonnées du registre.

Soient \(V=\RR^{12}\) et \(W=\RR^{\mathcal H}\), munis de leurs produits scalaires euclidiens. La matrice d'incidence définit l'application
\begin{equation}
 \mathcal I:V\longrightarrow W,
 \qquad (\mathcal I k)_H=\sum_{p\in H}k_p,
 \qquad \mathcal I_{H,p}=\mathbf1_{\{p\in H\}}.
 \label{rev4:ki:incidence-map}
\end{equation}
Le symbole \(\mathbf1\in V\) désigne le vecteur uniforme à douze
composantes. Nous définissons
\[
 \mu(k)=\frac1{12}\mathbf1^{\mathsf T}k,\qquad
 P_0=I_{12}-\frac1{12}\mathbf1\mathbf1^{\mathsf T},\qquad
 V_0=\mathbf1^\perp.
\]
Le projecteur \(P_0\) coïncide avec \(P_{11}\) de la section~\ref{sec:memory} ; l'indice zéro indique dans cette présentation la condition de moyenne nulle. La décomposition orthogonale est \(k=\mu(k)\mathbf1+P_0k\).

\begin{lemma}[Identité de Gram et normalisation incidentielle]
\label{rev4:ki:gram}
L’incidence satisfait
\begin{equation}
 \mathcal I^{\mathsf T}\mathcal I
       =36I_{12}+30\mathbf1\mathbf1^{\mathsf T}.
 \label{rev4:ki:gram-formula}
\end{equation}
Par conséquent,
\[
 U_W=\frac16\mathcal I\big|_{V_0}:
 V_0\longrightarrow E_{\rm inc}:=\mathcal I(V_0)
\]
est une isométrie sur un sous-espace de dimension onze.
\end{lemma}
\begin{proof}
Un point étant fixé, les quintuplets qui le contiennent sont au nombre de \(\binom{11}{4}\). Chaque hexade contenant ce point en comprend \(\binom54\). L'unicité de l'extension de chaque quintuplet donne \(66\) hexades par point. Une paire de points étant fixée, le même argument donne \(\binom{10}{3}/\binom43=30\) hexades par paire. L'entrée \((p,r)\) de \(\mathcal I^{\mathsf T}\mathcal I\) compte exactement les hexades contenant simultanément \(p\) et \(r\). Ses entrées diagonales sont donc \(66\), et les autres \(30\), ce qui prouve \eqref{rev4:ki:gram-formula}.

Pour \(z\in V_0\), le terme uniforme s’annule et
\[
 \left\|\frac16\mathcal I z\right\|^2
 =\frac1{36}z^{\mathsf T}
       \mathcal I^{\mathsf T}\mathcal I z
 =\|z\|^2.
\]
L’injectivité et la dimension de l’image se déduisent de cette
égalité. Le facteur six est ainsi fixé par le gramien du plan.
\end{proof}

Le dénombrement d’incidence détermine l’échelle de l’isométrie. Les
sommes sur les hexades sont des observables linéaires des positions ; leurs
produits scalaires contiennent la multiplicité avec laquelle chaque
position et chaque couple apparaissent dans le plan. L’élimination du
mode uniforme sépare ces deux niveaux de dénombrement et laisse le facteur
quadratique \(36\). Cette opération fournit un lien explicite
entre la représentation à douze positions et sa représentation
sur les hexades.

\subsection{Conservation de la moyenne et inversion des douze coordonnées}
\label{rev4:ki:full-subsection}

La composante uniforme est conservée au moyen d’une coordonnée scalaire
supplémentaire. Nous munissons \(\RR\oplus E_{\rm inc}\) du produit
scalaire
\[
 \langle(\mu,y),(\nu,z)\rangle_*=
 12\mu\nu+\langle y,z\rangle.
\]
Le poids douze est la norme au carré de \(\mathbf1\), fixée par le
nombre de positions. Nous définissons
\begin{equation}
 F:V\longrightarrow\RR\oplus E_{\rm inc},\qquad
 F(k)=\left(\mu(k),\frac16\mathcal I P_0k\right).
 \label{rev4:ki:F}
\end{equation}

\begin{theorem}[Isomorphisme isométrique du registre complet]
\label{rev4:ki:full-isometry}
L'application \(F\) est un isomorphisme linéaire isométrique. Son inverse est
\begin{equation}
 F^{-1}(\mu,y)=\mu\mathbf1+
                  \frac16P_0\mathcal I^{\mathsf T}y,
 \qquad y\in E_{\rm inc}.
 \label{rev4:ki:F-inverse}
\end{equation}
En particulier,
\begin{equation}
 12|\mu(k)|^2+
 \left\|\frac16\mathcal I P_0k\right\|^2=\|k\|^2.
 \label{rev4:ki:F-norm}
\end{equation}
\end{theorem}
\begin{proof}
La décomposition uniforme--centrée est orthogonale et le
lemme~\ref{rev4:ki:gram} conserve la norme de son second terme.
Cela démontre \eqref{rev4:ki:F-norm} et l’injectivité.
De plus,
\[
 \frac1{36}P_0\mathcal I^{\mathsf T}\mathcal I P_0=P_0.
\]
En substituant \(F(k)\) dans \eqref{rev4:ki:F-inverse}, on obtient \(\mu(k)\mathbf1+P_0k=k\). Réciproquement, si \(y\in E_{\rm inc}\), il existe un unique \(z\in V_0\) avec \(y=\mathcal I z/6\). La formule proposée renvoie \(\mu\mathbf1+z\), dont l'image est \((\mu,y)\). La surjectivité sur le codomaine déclaré et les deux identités d'inversion sont ainsi prouvées.
\end{proof}

La restriction \(y\in E_{\rm inc}\) exprime une compatibilité
linéaire effective. L’opérateur
\begin{equation}
 \Pi_{\rm inc}=\frac1{36}\mathcal I P_0
                         \mathcal I^{\mathsf T}:W\longrightarrow W
 \label{rev4:ki:inc-projector}
\end{equation}
est le projecteur orthogonal sur \(E_{\rm inc}\). En effet, il est
symétrique et son carré se réduit à lui-même au moyen de
\eqref{rev4:ki:gram-formula} ; il agit comme l’identité dans l’image
de \(\mathcal I P_0\). Par conséquent, la compatibilité se vérifie
au moyen de \(\Pi_{\rm inc}y=y\). Pour une donnée arbitraire de
\(W\), la reconstruction suivie de l’incidence restitue
\(\Pi_{\rm inc}y\) ; la composante orthogonale
\((I_W-\Pi_{\rm inc})y\) quantifie l’information externe à
cette représentation du registre.

\begin{proposition}[Covariance de la reconstruction]
\label{rev4:ki:equivariance}
Soit \(g\) un automorphisme du plan et soient \(\rho(g)\) et
\(\sigma(g)\) ses actions par permutation sur les positions et les
hexades. Alors
\[
 F\rho(g)=(1\oplus\sigma(g))F.
\]
La moyenne est conservée et la reconstruction commute avec le transport simultané des positions et de leurs hexades.
\end{proposition}
\begin{proof}
L’équivalence \(p\in H\Longleftrightarrow g(p)\in g(H)\)
implique \(\mathcal I\rho(g)=\sigma(g)\mathcal I\).
Les permutations fixent \(\mathbf1\), conservent la moyenne et
commutent avec \(P_0\). Substituer ces identités dans
\eqref{rev4:ki:F} démontre l’affirmation. La covariance de
l’inverse s’obtient en composant avec \(F^{-1}\).
\end{proof}

La normalisation isométrique et l'inversion arithmétique remplissent des fonctions distinctes. Dans la mise en coordonnées de Hadamard du registre, \(U=\mathcal H_{12}K\) et \(K=\mathcal H_{12}U/4\), avec \(\mathcal H_{12}^{\mathsf T}=\mathcal H_{12}\) et \(\mathcal H_{12}^2=4I_{12}\). L'opérateur \(J_H=\mathcal H_{12}/2\) est orthogonal. Par conséquent, \(FJ_H\) conserve la norme du registre signé normalisé \(U/2\) et le représente dans l'espace d'incidence élargi. L'inversion entière au moyen de \(\mathcal H_{12}/4\) conserve en outre le sous-réseau de congruences correspondant. Dans les deux cas, le registre engendré est transporté, avec l'ordre de phase fixé avant la transformation \cite{hmtX}.

\subsection{Composition exacte à partir des deux mémoires nonadiques}
\label{rev4:ki:memory-subsection}

Nous reprenons \(M\), \(q\), \(m_0\) et \(m_1\)
de \eqref{eq:mem-data}, avec l’opérateur de reconstruction
\(L=(M^*M|_{V_0})^{-1}M^*\) de la section~\ref{sec:memory}.
Le domaine compatible est
\(\mathcal D=\RR\times\operatorname{im}M\), avec
\(m=(m_0,m_1)\). Le théorème~\ref{thm:memory-inversion}
établit la reconstruction
\[
 \mathcal R(q,m)=\frac q{12}\mathbf1+Lm,
 \qquad LM=P_0.
\]
Comme \(L\) prend ses valeurs dans \(V_0\), la composition avec
l’incidence possède une expression directe :
\begin{equation}
 \mathcal C=F\mathcal R,\qquad
 \mathcal C(q,m)=\left(\frac q{12},\frac16\mathcal I Lm\right).
 \label{rev4:ki:memory-incidence}
\end{equation}

\begin{corollary}[Équivalence entre mémoire compatible et incidence intégrale]
\label{rev4:ki:composition}
L’application \(\mathcal C\) est un isomorphisme entre
\(\mathcal D\) et \(\RR\oplus E_{\rm inc}\), d’inverse
\begin{equation}
 \mathcal C^{-1}(\mu,y)=
 \left(12\mu,\frac16MP_0\mathcal I^{\mathsf T}y\right).
 \label{rev4:ki:C-inverse}
\end{equation}
Sa forme quadratique conservée est
\[
 \|\mathcal C(q,m)\|_*^2=
 \frac{|q|^2}{12}+\|Lm\|^2=\|\mathcal R(q,m)\|^2.
\]
\end{corollary}
\begin{proof}
La formule \eqref{rev4:ki:memory-incidence} résulte de
\(P_0\mathbf1=0\) et de \(P_0L=L\). L’inverse s’obtient
en appliquant \(k\mapsto(\mathbf1^{\mathsf T}k,Mk)\) à
\eqref{rev4:ki:F-inverse} ; on utilise \(M\mathbf1=0\).
Les identités de composition proviennent de \(LM=P_0\) et de
\(MLm=m\) pour \(m\in\operatorname{im}M\). L’égalité des
normes est \eqref{rev4:ki:F-norm} appliquée à \(\mathcal R(q,m)\).
\end{proof}

\begin{figure}[htbp]
\centering
\begin{tikzpicture}[>=Latex, every node/.style={font=\small},
  box/.style={draw,rounded corners=2pt,align=center,
              inner sep=6pt,minimum height=12mm}]
 \node[box] (memory) at (0,0) {Charge et mémoires\\$(q,m_0,m_1)$};
 \node[box] (register) at (4.1,0) {Registre intégral\\$K$};
 \node[box] (incidence) at (8.2,0) {Moyenne et incidence\\$(\mu,y)$};
 \node[box] (direction) at (4.1,-1.75)
    {Direction géométrique\\$\ell_K=\RR P_3K$};
 \draw[<->] (memory) -- node[above] {$\mathcal R$} (register);
 \draw[<->] (register) -- node[above] {$F$} (incidence);
 \draw[->] (register) -- node[right] {projection et passage à la droite}
    (direction);
\end{tikzpicture}
\caption{La charge et les deux mémoires compatibles reconstruisent les douze
coordonnées. La moyenne et la composante incidencielle fournissent une autre
représentation réversible du même registre. La direction géométrique
est une réduction ultérieure à fibres de dimension positive.}
\label{rev4:ki:figure}
\end{figure}

Cette composition détermine également la partition positive de \ref{sec:memory}. Si \(q>0\), ses longueurs se récupèrent au moyen de
\begin{equation}
 d_j=\frac1{12}+\frac{1}{6q}
             \bigl(P_0\mathcal I^{\mathsf T}y\bigr)_j,
 \qquad y=\frac16\mathcal I P_0K.
 \label{rev4:ki:positive-lengths}
\end{equation}
Les douze conditions \(d_j>0\) sont les équations de la coordinatisation positive exprimées en coordonnées d'incidence. Avec \(t_j=\sum_{r\leq j}d_r\), les mineurs ordonnés \(t_j-t_i=\sum_{i<r\leq j}d_r\) sont reconstruits à partir de \((\mu,y)\) ou de \((q,m_0,m_1)\). La positivité des mineurs et la conservation isométrique de l'incidence sont ainsi reliées par des applications explicites. La première est une inégalité dans la coordinatisation ordonnée ; la seconde est une identité de produits scalaires.

\subsection{Récupération de la polarisation et portée des réductions}
\label{rev4:ki:direction-subsection}

La réalisation des douze positions sur les classes de
\(A_5/C_5\) détermine le projecteur tridimensionnel
\[
 P_3=\frac3{60}\sum_{a\in A_5}
               \chi_3(a^{-1})\varrho(a).
\]
Ici, \(\varrho\) est la représentation par permutations et \(\chi_3\) le caractère tridimensionnel choisi dans la coordinatisation marquée. Ses représentants, valeurs de caractère et matrice effective sont fixés dans \eqref{eq:leech-projectors} ; le lemme~\ref{lem:leech-u} démontre \(P_3=P_3^{\mathsf T}=P_3^2\), \(\operatorname{rank}P_3=3\) et \(P_3\mathbf1=0\). Dans cette coordinatisation, la direction du registre s'obtient par
\begin{equation}
 u_K=P_3P_0K=P_3Lm
     =\frac16P_3P_0\mathcal I^{\mathsf T}y.
 \label{rev4:ki:direction-reconstruction}
\end{equation}
Les égalités découlent des deux inverses déjà démontrés. Ainsi,
la même direction se récupère à partir des mémoires ou de la
composante incidentielle ; la moyenne est conservée pour reconstruire
\(K\), même si le projecteur tridimensionnel l’annule.

Pour le registre de \eqref{eq:leech-K}, l’évaluation exacte
\eqref{eq:leech-u-norm} donne
\[
 \|u_K\|^2=
       \frac{6638585+2275584\sqrt5}{20}>0.
\]
La droite \(\ell_K=\RR u_K\) et le complément transversal
\(V_{10}(K)=V_0\cap u_K^\perp\) sont ainsi déterminés.
Son projecteur est
\begin{equation}
 P_{10}(K)=P_0-
       \frac{u_Ku_K^{\mathsf T}}{\|u_K\|^2}.
 \label{rev4:ki:rank-ten}
\end{equation}
En effet, le second terme est un projecteur orthogonal de rang un contenu dans \(V_0\) ; son produit avec \(P_0\), dans l'un ou l'autre ordre, coïncide avec lui-même. En élevant \eqref{rev4:ki:rank-ten} au carré, on récupère le même opérateur, et sa trace vaut \(11-1=10\). On obtient \(V_0=\ell_K\oplus^\perp V_{10}(K)\). Cette réduction est celle qui intervient ensuite dans la polarisation des douze plans réticulaires.

L’incidence transporte également ces opérateurs. Sur
\(E_{\rm inc}\), l’opérateur
\(\widehat P_3=U_WP_3U_W^*\) vérifie
\(\widehat P_3U_W=U_WP_3\) ; l’identité découle de
\(U_W^*U_W=I_{V_0}\). Transporter de la même manière
\(P_{10}(K)\) conserve son rang et son produit scalaire.
La correspondance entre les deux représentations s’établit
donc sur les vecteurs et les opérateurs, en plus de leurs dimensions.

La nécessité de conserver les différentes composantes admet
des vérifications précises. L’observation centrée de \(k\) coïncide
avec celle de \(k+c\mathbf1\) pour tout \(c\), tandis que leurs
charges diffèrent de \(12c\). En particulier, \(k=0\) et
\(k=\mathbf1\) ont la même composante incidentielle centrée
et des moyennes distinctes. La coordonnée \(\mu\) est exactement
l’information qui sépare ces deux registres. De même,
\(u_{ak+b\mathbf1}=a u_k\) pour \(a\ne0\) ; la droite
résultante coïncide avec \(\ell_k\), même si l’amplitude
et la moyenne changent. Ces identités décrivent les fibres des
applications sur l’espace ambiant déclaré.

L’extraction \(k\mapsto\{j:k_j\geq729\}\) possède d’autres
fibres. Pour le registre concret \(K\), aussi bien \(K\) que
\(K+e_1\) ont pour support \(\{4,6,9,10\}\), où
\(e_1\) est le premier vecteur de la base canonique. Ce sont deux vecteurs
ambiants différents ayant la même tétrade. Cette vérification
caractérise la perte d’information de l’application de support ; la
sélection terminale du registre reste fixée par sa
construction préalable. Le drapeau
\(B_K\subset H^-_\alpha\) conserve en outre l’hexade déterminée
par les signes antérieurs à la retenue, dont la coïncidence avec la
sélection incidencielle est démontrée dans la
proposition~\ref{prop:lf-flag}.

Le résultat commun de ces compositions est une continuité mathématique explicite. La mémoire compatible reconstruit le registre ; sa décomposition uniforme--centrée détermine une représentation incidentielle isométrique ; la reconstruction de ses longueurs détermine la coordinatisation positive ; et la projection tridimensionnelle détermine la direction de déformation. La tétrade et le drapeau conservent les incidences utilisées par les constructions exceptionnelles, tandis que \(F\) conserve la totalité du registre dodécaphasique. L'historique enrichi maintient à son tour l'orientation, le parcours et le prolongement qui situent ce registre dans la structure discrète conjointe du continu.
